\scanpage{108}\mspage{\handwritten{$8'$}}

The expression ``deduces from'' in B) means that we are adding arrows $h:K\handwritten{\to}F_i$ (each with preassigned source and target in $\underline{C}_n$), with as ``new axioms'' on \inserted{the bunch of} these uniquely the two relations $(\deleted{\unsure{a}}2)$ of \deleted{\unsure{B}}A), the category $\underline{C}_{n+1}$ being deduced from $\underline{C}_n$ in an obvious way, as the solution of a universal problem with\handwritten{in} the category of all categories where binary amalgamated products exist, -- and ``maps'' between these being functors which commut\unsure{e} to those fibered products.
% slanted
\margin{\handwritten{\unsure{unaccurate|inaccurate} -- amalgamated sums exist in $\underline{C}_{\unsure{o|0}}$ only, and \uline{these} should be \unsure{respected}\ldots}}
In practical terms, the arrows of $\underline{C}_{n+1}$ are those deduced from the arrows in $\underline{C}_n$ and the ``new'' arrows $h$, by combining formal operations of composing arrows by $v\handwritten{\circ}u$, and taking \inserted{\handwritten{(binary)}} amalgamated products of arrows. \deleted{\unsure{Of}}

NB Of course the condition $(1)$ in A) is necessary for the existence of an $h$ satisfying $(2)$. That it is sufficient too can be viewed as an extremely strong,\deleted{\unsure{version}} ``universal'' version of coherence conditions, concerning the various structures introduced on a $\handwritten{\infty}$-groupoid. Intuitively, it means that whenever we have \deleted{a}\inserted{two} way\handwritten{s} \deleted{to}of associat\unsure{ing} to a \inserted{finite} family \deleted{of} $(u_i)_{i\handwritten{\in}I}$ of objects of an $\handwritten{\infty}$-groupoid, $u_i\handwritten{\in}F_{n(i)}$, \deleted{\unsure{and such}} subjected to a ``standart'' set of relations on the $u_i$'s, an element of some $F_n$, in terms of the $\handwritten{\infty}$-groupoid structure only, then we have automatically \deleted{an isomorphism}\inserted{a ``homotopy''} between these built in in the very structure of the $\handwritten{\infty}$-groupoid, provided it makes at all a sense to ask for one (namely provided condition $(1)$ holds if $n\handwritten{\geq}1$). I have the feeling moreover that conditions (A) and (\unsure{B|3}) (plus the relation $\underline{C}_{\handwritten{\infty}}=\handwritten{\varinjlim}\underline{C}_n$) is all what will be ever needed, when using the definition of the structure species, -- plus of course the description of $\underline{C}_o$, and the \inserted{implicit} fact that the categories $\underline{C}_n$ are stable under binary fiber products and \inserted{\handwritten{that}} the inclusion functors commute to these.
% slanted
\margin{\handwritten{unaccurate -- see above!}}
Of course, the category which really interests us is $\underline{C}_{\handwritten{\infty}}$, the description of the intermediate $\underline{C}_n$\handwritten{'s} is merely technical \deleted{\unsure{but}} -- the main point is that there should exist an increasing sequence $(C_n)$ of subcategories of $\underline{C}_{\handwritten{\infty}}$, having same objects\handwritten{[}and ``same'' fiberproducts\handwritten{]}, such that $\underline{C}_{\handwritten{\infty}}$ should be the limit (i.e. every arrow in $C_{\handwritten{\infty}}$ should belong \deleted{\unsure{t}} to some $\underline{C}_n$), and such that the passage from $\underline{C}_n$ to $\underline{C}_{n+1}$ should satisfy (B). It is fairly obvious that these conditions alone do by no means caracterise $\underline{C}_{\handwritten{\infty}}$ up to equivalence, and still less the sequence of its subcategories $\underline{C}_n$. The point I wish to make though, before pursuing with a proposal of an explicit description, is that \uline{this ambiguity is in the nature of things}. \handwritten{R}oughly saying, two different mathematicians, working independently on the conceptual problem I had in mind, \deleted{\unsure{aski}}\unsure{assu}ming they both wind up with some explicit definition, will almost certainly get non equivalent definitions -- namely with non equivalent categories of (set-valued, say) $\handwritten{\infty}$-groupoids ! And, secondly and as importantly,

% archive sheet number: 107
